\section{Appendix: Rate comparison, branch formulas, and verification scope}

The support rule in \cref{obj:def:noise-support} completes the specification of the noisy alternatives. This appendix collects the comparisons that connect those alternatives and the public variance certificate to the resolution equation in \cref{obj:def:uniform-rate}. Three auxiliary results control changes in resolution, observable second moments, and cancellation information. Their roles in \cref{obj:thm:uniform-frontier} are followed by the branch interpretation and the larger-class converse.

\subsection{Noise cost and the two envelopes}

The noise cost summarizes the dependence of Gaussian inversion on endpoint resolution. Comparisons across resolutions must remain valid when the direct, intermediate, and compact branches meet. The following bounds provide that common control.

\begin{lemmav}{lem:noise-cost-scaling}[Noise cost scaling bounds]
For \(h\in(0,1]\) and \(\sigma\in[0,1]\), define the noise cost
\[
\Psi(h,\sigma)=
\begin{cases}
0, & \sigma\le h,\\
(\sigma/h)^{2/3}-1, & h<\sigma\ \text{and}\ \sigma^4\le h,\\
h^{-1/2}\bigl[1+\log(\sigma^2/\sqrt h)\bigr]-1,
& h<\sigma\ \text{and}\ h<\sigma^4.
\end{cases}
\]
For every \(\sigma\in[0,1]\), the function \(h\mapsto\Psi(h,\sigma)\) is continuous, nonnegative, and nonincreasing on \((0,1]\).

For all real \(h,t,\sigma\) satisfying \(0<h\le t\le\sigma\le1\),
\[
(t/h)^{1/2}
\le
\frac{1+\Psi(h,\sigma)}{1+\Psi(t,\sigma)}
\le t/h.
\]
For every \(h\in(0,1]\), \(\sigma\in[0,1]\), and real \(a\ge1\) satisfying \(ah\le1\),
\[
1+\Psi(ah,\sigma)
\le
\max\bigl\{1,\ a^{-1/2}[1+\Psi(h,\sigma)]\bigr\}.
\]
\end{lemmav}

Continuity keeps the cost consistent at the branch interfaces, while monotonicity orders finer and coarser endpoint resolutions. The ratio bounds quantify this ordering within the noisy region. The final inequality also accommodates a coarsening that reaches the direct region, where the cost vanishes. These properties support comparisons between the continuous resolution and the discrete polynomial choices used by the procedures and witnesses.

For the upper comparison, the relevant quantity is the integrated derivative sum entering the finite-sample radius. The next envelope expresses its dependence on degree and noise through the same cost function.

\begin{lemmav}{lem:variance-cost-envelope}[Uniform variance cost envelope]
Let \(m_L\) be the degree of the endpoint polynomial \(K_L\), and define its integrated derivative cost by
\[
V_L(\sigma)
=\frac34\sum_{j=0}^{m_L}\frac{\sigma^{2j}}{j!}
\int_0^1\bigl(K_L^{(j)}(x)\bigr)^2\,dx.
\]
For \(h>0\), define the noise cost by
\[
\Psi(h,\sigma)=
\begin{cases}
0, & \sigma\le h,\\
(\sigma/h)^{2/3}-1, & \sigma>h\ \text{and}\ \sigma^4\le h,\\
h^{-1/2}\bigl[1+\log(\sigma^2/\sqrt h)\bigr]-1,
& \sigma>h\ \text{and}\ \sigma^4>h.
\end{cases}
\]
There exists an absolute constant \(C>0\) such that, for every integer \(L\ge1\) and every \(\sigma\in[0,1]\),
\[
V_L(\sigma)\le C L^2
\exp\!\left\{C\bigl[1+\Psi(L^{-2},\sigma)\bigr]\right\}.
\]
\end{lemmav}

The envelope in \cref{obj:lem:variance-cost-envelope} places the derivative cost from \cref{obj:lem:exact-covariance} on the endpoint scale \(L^{-2}\). Together with the approximation bound in \cref{obj:lem:positive-endpoint}, it supplies the two components of the public radius in \cref{obj:synth_6}. Its range includes \(L=1\), whose polynomial degree is zero, and \(\sigma=0\), where the derivative sum retains its order-zero term.

The lower comparison uses the complete observed-law likelihood bound in \cref{obj:lem:full-marked-likelihood}. At the support and resolution specified by \cref{obj:def:noise-support}, its information bound takes a corresponding exponential form.

\begin{lemmav}{lem:cancellation-information-envelope}[Cancellation information envelope]
Let the parameters satisfy:
\begin{itemize}
\item \textbf{(Smoothness.)} \(0<\beta\le1\).
\item \textbf{(Noise scale.)} \(0<\sigma\le1\).
\item \textbf{(Polynomial degree.)} \(m\in\mathbb N\) and \(m\ge2\).
\end{itemize}
Write \(b=b_m(\sigma)\) and \(h=h_m(\sigma)\) as in
\cref{obj:def:noise-support}, and let \(P^+_{b,m}\) and \(P^-_{b,m}\)
be the alternative laws of \cref{obj:def:alternatives}, constructed using
\cref{obj:lem:classical-legendre-facts}. Set the perturbation constant
\(\kappa=1/100\), and define the noise cost by
\[
\Psi(h,\sigma)=
\begin{cases}
0, & \sigma\le h,\\
(\sigma/h)^{2/3}-1, & \sigma>h\ \text{and}\ \sigma^4\le h,\\
h^{-1/2}\bigl(1+\log(\sigma^2/\sqrt h)\bigr)-1,
& \sigma>h\ \text{and}\ \sigma^4>h.
\end{cases}
\]
For probability laws \(\mu,\nu\), use the extended chi-squared divergence
\[
\chi^2(\mu,\nu)=
\begin{cases}
\displaystyle\int\left(\frac{d\mu}{d\nu}-1\right)^2\,d\nu,
& \mu\ll\nu,\\
+\infty, & \mu\not\ll\nu.
\end{cases}
\]
Then the observed laws induced by the two alternatives at noise scale
\(\sigma\) satisfy
\[
\chi^2\bigl((P^+_{b,m})_{\mathrm{obs}},
           (P^-_{b,m})_{\mathrm{obs}}\bigr)
\le
\frac{32}{3}\kappa^2 h^{2\beta+1}
\exp\!\left(-\frac{1+\Psi(h,\sigma)}{32}\right).
\]
\end{lemmav}

The bound in \cref{obj:lem:cancellation-information-envelope} concerns the full observed triple, including the assignment and outcome marks. Its prefactor combines the squared perturbation amplitude with the endpoint resolution, while its exponential factor describes the information reduction associated with cancellation. Membership and target separation are supplied separately by \cref{obj:lem:legal-cancellation}. The degree range starts at \(m=2\), preserving the lowest admissible cancellation witness, and the positive-noise condition identifies the domain of this envelope.

\subsection{Uniform comparison and boundary cases}

The upper and lower comparisons in \cref{obj:thm:uniform-frontier} connect these envelopes to the balance equation. For the upper comparison, \cref{obj:thm:finite-certificate} already supplies measurable estimation, uniform coverage, and expected-length control by the selected public radius. The remaining rate comparison concerns that deterministic radius, using the approximation and variance bounds at integer polynomial indices.

The finite candidate set in \cref{obj:def:public-selection} is part of this comparison. Positive indices are restricted by the public cap, ties are resolved by the smallest minimizer, and index zero contributes the radius \(1/2\) specified in \cref{obj:synth_6}. Thus the upper bound in \cref{obj:thm:uniform-frontier} applies to the stated finite selection rule. The scaling control in \cref{obj:lem:noise-cost-scaling} provides the resolution comparisons needed for discrete indices, while the constant candidate supplies a valid certificate at coarse resolutions.

For the lower comparison at positive noise, \cref{obj:lem:cancellation-information-envelope} supplies observed-law proximity for the integer-degree alternatives of \cref{obj:def:noise-support}. The target separation and decision transfer in \cref{obj:lem:legal-cancellation,obj:lem:two-point-transfer} give the corresponding estimation-risk and honest-length bounds. The direct alternatives in \cref{obj:def:direct-alternatives} supply the direct-experiment component recorded in \cref{obj:thm:direct-reduction}. Both decision comparisons include the independent procedure seed.

Zero noise belongs to the exact direct branch of \cref{obj:thm:uniform-frontier}, with attainment and lower bounds also stated in \cref{obj:thm:direct-reduction}. At the direct interface, the resolution equals both the noise scale and \(n^{-1/(2\beta+1)}\). At the intermediate--compact interface, it equals \(\sigma^4\). Continuity in \cref{obj:lem:noise-cost-scaling} and the common constants in \cref{obj:thm:uniform-frontier} keep these boundary cases within the simultaneous comparison.

\subsection{Scalar branches and converse transfer}

The scalar coordinates in \cref{obj:thm:uniform-frontier} express the noisy resolution in two ways. On the intermediate branch, \(z=(\sigma/h_\star)^{2/3}\) measures the noise-to-resolution ratio. On the compact branch, \(\tau=h_\star^{-1/2}\) measures inverse square-root resolution. Their equations and admissible ranges specify the corresponding portions of the same balance equation. At the common boundary, both coordinates equal \(\sigma^{-2}\).

The sequence specializations in \cref{obj:thm:uniform-frontier} retain their stated quantifiers. For polynomially shrinking noise with a fixed exponent below \(1/(2\beta+1)\), the intermediate asymptotic order combines the noise power with a logarithmic improvement. At or above that exponent, the exact direct branch applies for every \(n\ge2\). Fixed positive noise eventually selects the compact branch. Uniformity of the original comparison also covers sequences that pass between branches, with constants independent of sample size and noise scale.

The larger-class implication is the target- and experiment-preserving embedding in \cref{obj:prop:published-class-converse-transfer}. The benchmark density and conditional-mean restrictions imply the local conditions of that class, and independence from the latent score and both potential outcomes implies the required joint independence from the score and observed outcome. Preservation of the entire randomized experiment makes the benchmark risk and honest-length lower bounds applicable to the larger decision problem. For positive noise, the published conditions invoked by this interpretation are Assumptions~1Y, 2C, and 7 of \citet{PeiShen2017DevilTails}.

\subsection{Direct-experiment specialization}

The direct alternatives and the zero-noise certificate identify the ordinary endpoint rate within the same benchmark experiment. The following result records a lower bound for the minimax absolute risk \leanref{sym:Risk}{\ensuremath{R_\beta(n,\sigma)}} and minimax honest expected length \leanref{sym:Length}{\ensuremath{\mathcal L_\beta(n,\sigma)}} at every admissible noise scale, together with attainment by the selected procedures at zero noise.

\begin{theoremv}{thm:direct-reduction}[Direct-experiment rates and attainment]
Fix \(0<\beta\le1\), where \(\beta\) is the Hölder exponent. Write \(R_\beta(n,\sigma)\) and \(\mathcal L_\beta(n,\sigma)\) for the real conversions of the extended minimax absolute risk and extended minimax honest expected length, respectively. Their extended objectives are
\[
\inf_{\widehat\theta}\sup_{P\in\mathcal M_\beta(\sigma)}
\mathbb E_{P,n}\!\left[|\widehat\theta-\theta(P)|\right],
\qquad
\inf_{I\in\mathcal I_{\beta,n,\sigma}}
\sup_{P\in\mathcal M_\beta(\sigma)}
\mathbb E_{P,n}\!\left[\operatorname{length}(I)\right].
\]
Here the first infimum ranges over admissible estimators, \(\mathcal M_\beta(\sigma)\) is the benchmark model class, \(\theta(P)\) is the cutoff treatment contrast, and \(\mathbb E_{P,n}\) denotes expectation under the joint law of the observed sample and independent procedure seed. The honest interval class \(\mathcal I_{\beta,n,\sigma}\) is defined in \cref{obj:def:honest-class}. Both minimizations are performed in the extended nonnegative reals before conversion to real numbers.

There exist constants \(c,C>0\), depending only on \(\beta\), such that, for every integer sample size \(n\ge2\) and every noise scale \(\sigma\in[0,1]\),
\[
R_\beta(n,\sigma)\ge c\,n^{-\beta/(2\beta+1)},
\qquad
\mathcal L_\beta(n,\sigma)\ge c\,n^{-\beta/(2\beta+1)}.
\]
At \(\sigma=0\),
\[
R_\beta(n,0)\le C\,n^{-\beta/(2\beta+1)},
\qquad
\mathcal L_\beta(n,0)\le C\,n^{-\beta/(2\beta+1)}.
\]
The selected estimator \(\widehat\theta_{L_\star}\) and selected interval \(I_{L_\star}\), evaluated at \((\beta,n,0)\), satisfy the extended worst-case bounds
\[
\sup_{P\in\mathcal M_\beta(0)}
\mathbb E_{P,n}\!\left[
|\widehat\theta_{L_\star}-\theta(P)|
\right]
\le C\,n^{-\beta/(2\beta+1)},
\qquad
\sup_{P\in\mathcal M_\beta(0)}
\mathbb E_{P,n}\!\left[
\operatorname{length}(I_{L_\star})
\right]
\le C\,n^{-\beta/(2\beta+1)}.
\]
Moreover, \(I_{L_\star}\in\mathcal I_{\beta,n,0}\), the honest interval class of \cref{obj:def:honest-class}.
\end{theoremv}

At zero noise, the endpoint approximation and the sampling variation in the public certificate recover the Hölder endpoint rate. The matching lower bounds identify that rate as the benchmark difficulty for both estimation and honest connected intervals. Their validity throughout \(\sigma\in[0,1]\) supplies the direct-experiment baseline for the uniform comparison, while the zero-noise attainment clauses connect that baseline to the observable procedures already specified in \cref{obj:def:public-selection}.

\subsection{Precision summary}

The established relationships can now be collected around the causal target, the public certificate, and the two families of alternatives. The Lean component definitions are total on their ambient natural-number and real domains, including the index-zero fallback and values outside the statistical domain. The guarantees below apply only for \(n\ge2\) and \(\sigma\in[0,1]\), while positive kernel indices and cancellation degrees use the ranges stated in their defining environments. The following object records these components with those domain distinctions and their defining references.

\begin{remarkv}{def:frontier-handle}[Endpoint precision handle]
The endpoint precision handle comprises the following objects, for \(0<\beta\le1\), \(n\in\mathbb N\), and \(\sigma\in\mathbb R\), given the classical Legendre identities in \cref{obj:lem:classical-legendre-facts}.
\begin{itemize}
\item \textbf{(Target.)} The cutoff causal contrast
\[
\theta(P)=\mu_1(P,0)-\mu_0(P,0).
\]
\item \textbf{(Variance.)} For every \(L\in\mathbb N\), the finite derivative sum
\[
m_L=\deg K_L,\qquad
V_L(\sigma)=\frac34\sum_{j=0}^{m_L}
\frac{\sigma^{2j}}{j!}
\int_0^1\bigl(K_L^{(j)}(x)\bigr)^2\,dx,
\]
where \(K_L\) is defined in \cref{obj:def:endpoint-kernel}.
\item \textbf{(Certificate.)} The selected public radius
\[
E_\beta(n,\sigma)=\rho_{L_\star},
\]
with the radius and selected index specified in \cref{obj:def:public-selection}.
\item \textbf{(Cancellation witnesses.)} For every integer \(m\ge2\), when \(0<\sigma\le1\), and for either sign, the alternative laws
\[
b_m(\sigma)=\min\left\{1,\frac{\sigma\sqrt m}{8}\right\},
\qquad
P^\pm_{b_m(\sigma),m},
\]
defined in \cref{obj:def:alternatives,obj:def:noise-support}.
\item \textbf{(Direct witnesses.)} For every \(0<h\le1\) and either sign, the alternative laws \(P_h^{\mathrm{dir},\pm}\) defined in \cref{obj:def:direct-alternatives}.
\item \textbf{(Rate and procedures.)} The resolution rate, selected measurable estimator, and selected interval procedure
\[
r_\beta(n,\sigma),\qquad
\widetilde\theta_{\beta,n,\sigma},\qquad
\widetilde I_{\beta,n,\sigma},
\]
defined in \cref{obj:def:uniform-rate}.
\end{itemize}
\end{remarkv}

The precision handle in \cref{obj:def:frontier-handle} collects quantities already defined and used above. The next proposition consolidates their relationships, including the attainable public radius, the minimax comparisons, and the scalar branch equations of \cref{obj:thm:uniform-frontier}.

\begin{propositionv}{thm:uniform-frontier-resolution}[Uniform precision frontier]
Fix \(\beta\in(0,1]\), and write \(\nu=2\beta+1\). For \(n\ge2\) and \(\sigma\in[0,1]\), put \(h_0=n^{-1/\nu}\) and define
\[
\Psi(h,\sigma)=
\begin{cases}
0, & \sigma\le h,\\
(\sigma/h)^{2/3}-1, & \sigma^4\le h<\sigma,\\
h^{-1/2}\bigl\{1+\log(\sigma^2h^{-1/2})\bigr\}-1,
    & 0<h<\sigma^4.
\end{cases}
\]
There is a unique \(h_\star(n,\sigma)\in[h_0,1]\) satisfying
\[
\nu\log\!\left(\frac1{h_\star(n,\sigma)}\right)
+\Psi\bigl(h_\star(n,\sigma),\sigma\bigr)=\log n.
\]
Set
\[
r_\beta(n,\sigma)=h_\star(n,\sigma)^\beta,\qquad
\widetilde\theta_{\beta,n,\sigma}=\widehat\theta_{L_\star},
\qquad
\widetilde I_{\beta,n,\sigma}=I_{L_\star}.
\]

There exist constants \(c,C>0\), depending on \(\beta\) and independent of \(n\) and \(\sigma\), such that, simultaneously for every integer \(n\ge2\) and every \(\sigma\in[0,1]\),
\[
\begin{aligned}
c\,r_\beta(n,\sigma)
&\le R_\beta(n,\sigma)
\le \sup_{P\in\mathcal M_\beta(\sigma)}
   \mathbb E_{P,n}
   \left|\widetilde\theta_{\beta,n,\sigma}-\theta(P)\right|
\le E_\beta(n,\sigma)
\le C\,r_\beta(n,\sigma),\\
c\,r_\beta(n,\sigma)
&\le \mathcal L_\beta(n,\sigma)
\le \sup_{P\in\mathcal M_\beta(\sigma)}
   \mathbb E_{P,n}
   \operatorname{length}\bigl(\widetilde I_{\beta,n,\sigma}\bigr)
\le 2E_\beta(n,\sigma)
\le C\,r_\beta(n,\sigma),\\
\widetilde I_{\beta,n,\sigma}
&\in\mathcal I_{\beta,n,\sigma}.
\end{aligned}
\]
The estimation and interval comparisons range over the declared measurable estimator and honest connected-interval decision classes, including procedures with an independent uniform random seed.

The resolution has the following branches:
\begin{itemize}
\item \textbf{(Direct branch.)}
If \(\sigma\le n^{-1/\nu}\), then
\[
h_\star(n,\sigma)=n^{-1/\nu}.
\]

\item \textbf{(Intermediate branch.)}
If \(\sigma>n^{-1/\nu}\) and
\[
\log(n\sigma^{4\nu})\le \sigma^{-2}-1,
\]
then
\[
h_\star(n,\sigma)=\sigma z^{-3/2},
\qquad
z+\frac{3\nu}{2}\log z=1+\log(n\sigma^\nu),
\qquad
1<z\le\sigma^{-2},
\]
where the displayed equation uniquely determines \(z\).

\item \textbf{(Compact branch.)}
If \(\sigma>n^{-1/\nu}\) and
\[
\log(n\sigma^{4\nu})\ge \sigma^{-2}-1,
\]
then
\[
h_\star(n,\sigma)=\tau^{-2},
\qquad
2\nu\log\tau+
\tau\bigl\{1+\log(\sigma^2\tau)\bigr\}=\log n+1,
\qquad
\tau\ge\sigma^{-2},
\]
where the displayed equation uniquely determines \(\tau\).
\end{itemize}
The intermediate and compact equations agree at
\(h_\star=\sigma^4\); the direct interface is
\(h_\star=\sigma=n^{-1/\nu}\). The bounds use the same constants along every sequence \(\sigma=\sigma_n\in[0,1]\), including shrinking-noise sequences.

In particular, for \(\beta=1\) and \(\sigma=n^{-1/6}\), the intermediate branch applies for every \(n\ge2\), and
\[
r_1(n,n^{-1/6})
\asymp
\frac{n^{-1/6}}{\bigl[1+\tfrac12\log n\bigr]^{3/2}},
\]
with comparison constants independent of \(n\). Throughout, the observed side label is \(D=\mathbf1\{X\ge0\}\).
\end{propositionv}

The comparisons in \cref{obj:thm:uniform-frontier-resolution} place minimax risk, honest expected length, and the selected public radius on the same resolution scale. The variance and cancellation envelopes connect attainable precision and statistical indistinguishability through the common noise cost. Because the constants apply simultaneously over the stated sample sizes and noise scales, the branch formulas describe both fixed-noise and shrinking-noise experiments with the same comparisons.

\subsection{Verification scope}

The statements and proofs are machine-checked in Lean 4, including the finite-sample certificate, the auxiliary bounds, the uniform and direct-experiment comparisons, and the converse transfer. Their statistical scope uses the Gaussian calibration, independence, bounded-support density restrictions, conditional-mean bounds, and full-interval Hölder restrictions in \cref{obj:def:model}, together with the declared sampling and decision classes. Lean definitions are total on their ambient types; statistical guarantees use the stated domains \(n\ge2\), \(\sigma\in[0,1]\), and the specified positive-index ranges, with index-zero procedures defined separately. The representation uses finite polynomial inverse-heat sums, induced observed and product laws, and extended nonnegative minimax objectives before real conversion. The classical Legendre and Hermite identities are established by verified declarations, and the supplied theorem-local verification metadata records empty external proof-dependency lists. The citation to \citet[Assumptions~1Y, 2C, and 7]{PeiShen2017DevilTails} supplies the published conditions used for the larger-class interpretation.
